% This section lists a number of 2-3 research questions that the thesis should address. 
% The research questions should be put into the context of the idea, and you should explain 
% what they mean and how you want to address each. We also suggest including an example scenario 
% of a successful outcome of the thesis, i.e., what would your tool allow us to do once it exists?

This paper will address the following research questions:
\begin{enumerate}
    \item How can automated model checking and other formal-verification techniques be used for verifying liveness (replication reliability) and safety properties 
        in self-replicating robotic systems(SRRS), and what domain-specific adaptations are required for practical deployment?
    \item What systematic methodology can translate the behaviours and structural constraints 
        of self-replicating robots into expressive temporal-logic specifications?
    \item What fundamental limitations—computational, semantic, or modelling-related—do current 
         face when applied to highly dynamic, self-replicating robotic collectives, and how might these limitations be mitigated?
\end{enumerate}

\section{Approach}

\subsubsection{Experimental platform}
As a subject of a research a simple robot will be taken, with two grippers and an elbow joint that can manipulate blocks (cuboctahedral voxels) one at a time.
An experimental platform will be realised as a simulation in Gazebo environment with control from ROS2 nodes. The  blocks and environment geometry robot joints and sensors and conections 
will be defined in Xacro/URDF files. 


\subsubsection{Generation of a model and specifications for formal-verification}
A custom program has to be written in Python, that parses the Xacro files and ROS2 Nodes controller to generate a NuSMV model:

  - a Boolean occupancy matrix representing the neighbourhood on a field relevant to the current replication.

  - replication program-counter plus local flags (gripper-occupied, block -held, etc.).

  The liveness and safety properties has to be defined independent from the simulation, in a form close to the following:
  \begin{itemize}
    
      \item Liveness. $G (start \rightarrow F (Robot_1==ready \land Robot_2==ready))$ — once replication starts, eventually both robots are able to function independently.

      \item Collision-free safety. $G \neg collision$ (or $AG \neg collision$ for CTL).

\end{itemize}


% \section{Analysis of results}


% \section{Evaluation}

